\documentclass[aspectratio=169]{beamer}
\usetheme{Madrid}
\usecolortheme{default}
\usepackage[utf8]{inputenc}
\usepackage[spanish]{babel}
\usepackage{graphicx}
\usepackage{hyperref}

\title{Normas APA - Séptima Edición}
\subtitle{Formato General del Documento}
\author{Edison Achalma}
\institute{Universidad Nacional de San Cristóbal de Huamanga}
\date{\today}

\begin{document}

\frame{\titlepage}

\begin{frame}{Tabla de Contenido}
\tableofcontents
\end{frame}

\section{Configuración de Página}

\begin{frame}{Tamaño del Papel}
\begin{block}{Especificaciones}
\begin{itemize}
    \item Tamaño: Carta (Letter)
    \item Dimensiones: 21.59 cm × 27.94 cm (8.5" × 11")
\end{itemize}
\end{block}

\vspace{0.5cm}
% Imagen: Representación visual del tamaño carta con dimensiones
\begin{center}
\includegraphics[width=0.5\textwidth]{images/tamano_papel.png}
\end{center}
\end{frame}

\begin{frame}{Márgenes del Documento}
% Imagen: Esquema de página mostrando los 4 márgenes de 2.54 cm
\begin{center}
\includegraphics[width=0.7\textwidth]{images/margenes_documento.png}
\end{center}

\begin{block}{Todos los márgenes}
2.54 cm (1 pulgada) en los cuatro lados
\end{block}
\end{frame}

\section{Fuentes Tipográficas}

\begin{frame}{Fuentes Permitidas}
\textbf{APA 7 permite varias fuentes:}

\vspace{0.3cm}
\textbf{Fuentes CON serifa:}
\begin{itemize}
    \item Times New Roman - 12 puntos
    \item Georgia - 11 puntos
    \item Computer Modern - 10 puntos
\end{itemize}

\vspace{0.3cm}
\textbf{Fuentes SIN serifa:}
\begin{itemize}
    \item Calibri - 11 puntos
    \item Arial - 11 puntos
    \item Lucida Sans Unicode - 10 puntos
\end{itemize}
\end{frame}

\begin{frame}{Ejemplos de Fuentes}
% Imagen: Muestra del mismo texto en las diferentes fuentes permitidas
\begin{center}
\includegraphics[width=0.85\textwidth]{images/ejemplos_fuentes.png}
\end{center}

\begin{alertblock}{Consistencia}
Use la misma fuente en todo el documento
\end{alertblock}
\end{frame}

\begin{frame}{Diferencia: Con serifa vs Sin serifa}
% Imagen: Comparación visual entre fuentes con serifa y sin serifa, mostrando las "patitas"
\begin{center}
\includegraphics[width=0.8\textwidth]{images/serifa_vs_sin_serifa.png}
\end{center}
\end{frame}

\section{Espaciado}

\begin{frame}{Interlineado}
\begin{block}{Regla General}
\textbf{Interlineado doble (2.0)} en todo el documento
\end{block}

\textbf{Incluye:}
\begin{itemize}
    \item Texto principal
    \item Referencias
    \item Notas al pie
    \item Títulos y subtítulos
\end{itemize}

\vspace{0.3cm}
% Imagen: Ejemplo de texto con interlineado 2.0 vs otros espaciados
\begin{center}
\includegraphics[width=0.6\textwidth]{images/ejemplo_interlineado.png}
\end{center}
\end{frame}

\begin{frame}{Espaciado entre Párrafos}
\begin{alertblock}{Importante}
\textbf{NO} se debe dejar espacio adicional entre párrafos
\end{alertblock}

% Imagen: Comparación de texto CON espacio entre párrafos (incorrecto) y SIN espacio (correcto)
\begin{center}
\includegraphics[width=0.75\textwidth]{images/espaciado_parrafos.png}
\end{center}
\end{frame}

\section{Alineación}

\begin{frame}{Alineación del Texto}
\begin{block}{Regla}
Alineación a la \textbf{izquierda}, sin justificar
\end{block}

\vspace{0.5cm}
% Imagen: Comparación de texto alineado a la izquierda vs justificado
\begin{center}
\includegraphics[width=0.8\textwidth]{images/alineacion_texto.png}
\end{center}
\end{frame}

\begin{frame}{¿Por qué NO justificar?}
\textbf{Razones para usar alineación izquierda:}
\begin{itemize}
    \item Evita espacios irregulares entre palabras
    \item Mejora la legibilidad
    \item Es el estándar de APA
    \item Facilita la lectura en pantalla
\end{itemize}

\vspace{0.3cm}
% Imagen: Ejemplo de problemas de espaciado en texto justificado
\begin{center}
\includegraphics[width=0.65\textwidth]{images/problemas_justificado.png}
\end{center}
\end{frame}

\section{Sangría}

\begin{frame}{Sangría en Párrafos}
\begin{block}{Primera línea de cada párrafo}
Sangría de \textbf{1.27 cm} (0.5 pulgadas)
\end{block}

% Imagen: Ejemplo de párrafo con sangría de 1.27 cm marcada
\begin{center}
\includegraphics[width=0.75\textwidth]{images/sangria_parrafo.png}
\end{center}
\end{frame}

\begin{frame}{Sangría Francesa}
\textbf{Se usa únicamente en las Referencias}

\vspace{0.3cm}
% Imagen: Ejemplo de sangría francesa en lista de referencias
\begin{center}
\includegraphics[width=0.8\textwidth]{images/sangria_francesa.png}
\end{center}

\begin{block}{Formato}
Primera línea sin sangría, las siguientes con sangría de 1.27 cm
\end{block}
\end{frame}

\begin{frame}{Tipos de Sangría}
% Imagen: Comparación visual de sangría normal vs sangría francesa
\begin{center}
\includegraphics[width=0.9\textwidth]{images/tipos_sangria.png}
\end{center}
\end{frame}

\section{Numeración de Páginas}

\begin{frame}{Números de Página}
\begin{block}{Ubicación}
Esquina superior derecha de cada página
\end{block}

\textbf{Características:}
\begin{itemize}
    \item Números arábigos (1, 2, 3...)
    \item Desde la portada
    \item En todas las páginas
    \item Sin palabra "página" o "pág."
\end{itemize}

\vspace{0.3cm}
% Imagen: Ejemplo de página con número en esquina superior derecha
\begin{center}
\includegraphics[width=0.5\textwidth]{images/numeracion_pagina.png}
\end{center}
\end{frame}

\section{Formato de Títulos}

\begin{frame}{Formato de los 5 Niveles}
% Imagen: Tabla completa con los 5 niveles de títulos y su formato específico
\begin{center}
\includegraphics[width=0.95\textwidth]{images/formato_5_niveles.png}
\end{center}
\end{frame}

\begin{frame}{Niveles 1, 2 y 3}
% Imagen: Ejemplo práctico mostrando niveles 1, 2 y 3 en un texto
\begin{center}
\includegraphics[width=0.8\textwidth]{images/ejemplo_niveles_1_2_3.png}
\end{center}
\end{frame}

\begin{frame}{Niveles 4 y 5}
% Imagen: Ejemplo práctico mostrando niveles 4 y 5 (con sangría, texto en la misma línea)
\begin{center}
\includegraphics[width=0.8\textwidth]{images/ejemplo_niveles_4_5.png}
\end{center}
\end{frame}

\section{Abreviaturas}

\begin{frame}{Abreviaturas Estándar de APA}
% Imagen: Tabla completa de abreviaturas (cap., ed., trad., p., pp., Vol., etc.)
\begin{center}
\includegraphics[width=0.85\textwidth]{images/tabla_abreviaturas.png}
\end{center}
\end{frame}

\begin{frame}{Uso de Abreviaturas}
\begin{block}{Reglas}
\begin{itemize}
    \item Usar solo abreviaturas estándar de APA
    \item Primera vez: término completo [abreviatura]
    \item Siguientes veces: solo abreviatura
    \item Evitar abreviaturas en títulos
\end{itemize}
\end{block}

% Imagen: Ejemplo de introducción de abreviatura en texto
\begin{center}
\includegraphics[width=0.7\textwidth]{images/ejemplo_uso_abreviaturas.png}
\end{center}
\end{frame}

\section{Listas}

\begin{frame}{Formato de Listas}
\textbf{Tipos de listas:}
\begin{itemize}
    \item Listas con viñetas (bullet points)
    \item Listas numeradas
    \item Listas dentro del texto
\end{itemize}

\vspace{0.5cm}
% Imagen: Ejemplos de los tres tipos de listas correctamente formateadas
\begin{center}
\includegraphics[width=0.75\textwidth]{images/tipos_listas.png}
\end{center}
\end{frame}

\begin{frame}{Cuándo Usar Listas}
\begin{alertblock}{APA recomienda}
Evitar el uso excesivo de listas. Preferir la redacción en prosa cuando sea posible.
\end{alertblock}

\textbf{Usar listas cuando:}
\begin{itemize}
    \item Se presentan 3+ elementos relacionados
    \item Se necesita claridad visual
    \item Son pasos de un procedimiento
    \item El usuario lo solicita explícitamente
\end{itemize}
\end{frame}

\section{Resumen del Formato}

\begin{frame}{Checklist de Formato}
% Imagen: Lista de verificación visual con todos los elementos de formato
\begin{center}
\includegraphics[width=0.85\textwidth]{images/checklist_formato.png}
\end{center}
\end{frame}

\begin{frame}{Configuración Rápida en Word}
% Imagen: Capturas de pantalla mostrando dónde configurar cada elemento en MS Word
\begin{center}
\includegraphics[width=0.9\textwidth]{images/configuracion_word.png}
\end{center}
\end{frame}

\begin{frame}[standout]
\Huge ¿Preguntas?
\end{frame}

\begin{frame}{Próxima Clase}
\begin{center}
\Large \textbf{Estilo de Redacción Académica}
\vspace{1cm}

Veremos:
\begin{itemize}
    \item Lenguaje formal y objetivo
    \item Voz activa vs pasiva
    \item Lenguaje inclusivo
    \item Evitar sesgos
\end{itemize}
\end{center}
\end{frame}

\end{document}
